\section{Modelado de Casos de Uso y Procesos Operativos (UML)}
\label{sec:casos_uso_actividades}
% =============================================================================

\subsection{Diagrama de Casos de Uso UML}

La Figura \ref{fig:uml_casos_uso} expone los Casos de Uso del sistema, detallando los actores involucrados, los límites de la plataforma y las relaciones de inclusión (\guillemotleft include\guillemotright) y extensión (\guillemotleft extend\guillemotright).

\begin{figure}[htbp]
	\centering
	\begin{tikzpicture}[
		actor/.style={rectangle, draw=PQCMainBlue, fill=BoxBgBlue, thick, rounded corners=3pt, minimum width=2.4cm, minimum height=0.8cm, font=\bfseries\scriptsize, align=center},
		uc/.style={ellipse, draw=PQCDark, fill=white, thick, minimum width=4.0cm, minimum height=0.7cm, font=\scriptsize, align=center},
		systembox/.style={rectangle, draw=gray!60, dashed, fill=gray!3, rounded corners=6pt, inner sep=10pt},
		assoc/.style={thick, draw=black!75},
		include/.style={-{Latex[length=2mm]}, dashed, draw=PQCTeal, font=\tiny\bfseries},
		extend/.style={-{Latex[length=2mm]}, dashed, draw=PQCCrimson, font=\tiny\bfseries}
	]
		% Actores Izquierda
		\node[actor] (votante) at (-5.4, 0.8) {Votante\\(Ciudadano)};
		\node[actor] (ar) at (-5.4, -4.5) {Autoridad de\\Registro};
		
		% Actores Derecha
		\node[actor] (sv) at (5.4, 0.8) {Servidor de\\Votación};
		\node[actor] (ac) at (5.4, -4.5) {Autoridad de\\Conteo};
		
		% Casos de Uso
		\node[uc] (uc1) at (0, 3.0) {CU1: Registrarse en Padrón};
		\node[uc] (uc2) at (0, 1.8) {CU2: Emitir Credencial de Votación};
		\node[uc] (uc3) at (0, 0.6) {CU3: Cifrar Voto en Cabina};
		\node[uc] (uc4) at (0, -0.6) {CU4: Firmar Boleta Digitalmente};
		\node[uc] (uc5) at (0, -1.8) {CU5: Validar Boleta y Duplicados};
		\node[uc] (uc6) at (0, -3.0) {CU6: Custodiar en Ánfora y Recibo};
		\node[uc] (uc7) at (0, -4.2) {CU7: Descifrar Urna en Conteo};
		\node[uc] (uc8) at (0, -5.4) {CU8: Escrutinio y Auditoría};
		\node[uc, draw=PQCCrimson] (uc9) at (0, -6.6) {CU9: Alerta Intento Doble Voto};
		
		% Límite del Sistema
		\begin{scope}[on background layer]
			\node[systembox, fit=(uc1) (uc9), label={[font=\bfseries\small, text=PQCDark]north:Límite del Sistema de Votación Post-Cuántica (FIPS 203 / 204)}] (sys) {};
		\end{scope}
		
		% Relaciones Votante
		\draw[assoc] (votante) -- (uc1);
		\draw[assoc] (votante) -- (uc3);
		
		% Relaciones AR
		\draw[assoc] (ar) -- (uc1);
		\draw[assoc] (ar) -- (uc2);
		
		% Relaciones SV
		\draw[assoc] (sv) -- (uc5);
		\draw[assoc] (sv) -- (uc6);
		
		% Relaciones AC
		\draw[assoc] (ac) -- (uc7);
		\draw[assoc] (ac) -- (uc8);
		
		% Inclusiones y extensiones
		\draw[include] (uc3) -- (uc4) node[midway, right] {\guillemotleft include\guillemotright};
		\draw[include] (uc5) -- (uc6) node[midway, right] {\guillemotleft include\guillemotright};
		\draw[include] (uc7) -- (uc8) node[midway, right] {\guillemotleft include\guillemotright};
		\draw[extend] (uc9) -- (uc5) node[midway, right] {\guillemotleft extend\guillemotright};
	\end{tikzpicture}
	\caption{Diagrama de Casos de Uso UML del sistema.}
	\label{fig:uml_casos_uso}
	\par\vspace{3pt}{\footnotesize\textit{Fuente: Elaboración propia.}}
\end{figure}

El diagrama formaliza la participación de los actores:
\begin{itemize}
	\item \textbf{CU1 y CU2 (Empadronamiento):} El votante se acredita de forma presencial o institucional ante la AR, quien genera y firma digitalmente su credencial ML-DSA-65.
	\item \textbf{CU3 y CU4 (Emisión de Boleta):} En el dispositivo del votante, formular una opción electoral incluye de manera forzosa el encapsulamiento simétrico ML-KEM-768 y la firma digital de la boleta con la clave privada del votante.
	\item \textbf{CU5 y CU6 (Custodia y Recibo):} El Servidor valida la credencial de AR, comprueba que el votante no figure con sufragio previo y almacena la boleta en la urna. En caso de detectarse un intento de duplicación, se dispara el caso de uso extendido CU9, abortando la transacción y registrando la alerta en auditoría.
	\item \textbf{CU7 y CU8 (Escrutinio Oficial):} Tras la clausura comicial, la Autoridad de Conteo utiliza su clave privada para desencapsular en bloque las claves simétricas, descifrar los votos, tabular los resultados y emitir la lista pública de comprobantes SHA3-256.
\end{itemize}

\subsection{Diagrama de Actividades UML: Flujo Operativo de Extremo a Extremo}

La Figura \ref{fig:uml_actividades} detalla el Diagrama de Actividades UML, ilustrando el flujo secuencial, las decisiones lógicas y las bifurcaciones condicionales que se ejecutan durante la emisión del voto.

\begin{figure}[htbp]
	\centering
	\begin{tikzpicture}[
		scale=0.9,
		act/.style={rectangle, draw=PQCMainBlue, fill=BoxBgBlue, thick, rounded corners=5pt, font=\scriptsize\sffamily, align=center, minimum height=0.75cm, text width=3.8cm},
		dec/.style={diamond, draw=PQCAmber, fill=BoxBgAmber, aspect=2.0, font=\tiny\bfseries, align=center, inner sep=2pt},
		flow/.style={-{Latex[length=2.5mm]}, thick, draw=black!85},
		init/.style={circle, fill=black, inner sep=4.5pt},
		stop/.style={circle, draw=black, thick, fill=white, inner sep=4.5pt}
	]
		% Nodo Inicial
		\node[init] (start) at (0, 0) {};
		
		% Actividades Cliente
		\node[act, below=0.6cm of start] (a1) {1. Presentar credencial y seleccionar candidato};
		\node[act, below=0.6cm of a1] (a2) {2. Generar clave secreta temporal y protegerla\\con la clave pública de la Autoridad de Conteo};
		\node[act, below=0.6cm of a2] (a3) {3. Cifrar el contenido del voto con AES-256\\y firmar digitalmente con la clave del votante};
		\node[act, below=0.6cm of a3] (a4) {4. Transmitir paquete de voto cifrado\\al Servidor por conexión segura};
		
		% Decisiones Servidor
		\node[dec, below=0.7cm of a4] (d1) {¿Credencial de AR\\válida y en padrón?};
		\node[dec, below=0.9cm of d1] (d2) {¿Firma digital del\\votante es auténtica?};
		\node[dec, below=0.9cm of d2] (d3) {¿Votante registra\\sufragio previo?};
		
		% Rechazos
		\node[act, fill=red!10, draw=PQCCrimson, right=2.0cm of d1, text width=2.8cm] (err1) {Rechazar: Credencial no autorizada};
		\node[act, fill=red!10, draw=PQCCrimson, right=2.0cm of d2, text width=2.8cm] (err2) {Rechazar: Firma digital corrupta};
		\node[act, fill=red!10, draw=PQCCrimson, right=2.0cm of d3, text width=2.8cm] (err3) {Rechazar: Intento de voto duplicado};
		
		% Actividades Éxito
		\node[act, fill=BoxBgGreen, draw=PQCGreen, below=0.9cm of d3] (a5) {5. Marcar elector como 'Ya Votó'\\y archivar boleta anónima en ánfora};
		\node[act, fill=BoxBgGreen, draw=PQCGreen, below=0.6cm of a5] (a6) {6. Generar código de verificación (SHA-3)\\y devolver recibo firmado por el Servidor};
		\node[act, below=0.6cm of a6] (a7) {7. Votante descarga/imprime recibo\\con código QR para auditoría};
		
		% Nodo Final
		\node[stop, below=0.6cm of a7] (final) {};
		\node[circle, fill=black, inner sep=2.5pt] at (final.center) {};
		
		% Conectores
		\draw[flow] (start) -- (a1);
		\draw[flow] (a1) -- (a2);
		\draw[flow] (a2) -- (a3);
		\draw[flow] (a3) -- (a4);
		\draw[flow] (a4) -- (d1);
		
		\draw[flow] (d1) -- node[left, font=\tiny\bfseries] {Sí} (d2);
		\draw[flow] (d1) -- node[above, font=\tiny\bfseries] {No} (err1);
		
		\draw[flow] (d2) -- node[left, font=\tiny\bfseries] {Sí} (d3);
		\draw[flow] (d2) -- node[above, font=\tiny\bfseries] {No} (err2);
		
		\draw[flow] (d3) -- node[left, font=\tiny\bfseries] {No} (a5);
		\draw[flow] (d3) -- node[above, font=\tiny\bfseries] {Sí} (err3);
		
		\draw[flow] (a5) -- (a6);
		\draw[flow] (a6) -- (a7);
		\draw[flow] (a7) -- (final);
		
		% Retornos de error al final
		\node[stop, right=5.4cm of d2] (err_final) {};
		\node[circle, fill=black, inner sep=2.5pt] at (err_final.center) {};
		\draw[flow, draw=PQCCrimson] (err1) -| (err_final);
		\draw[flow, draw=PQCCrimson] (err2) -- (err_final);
		\draw[flow, draw=PQCCrimson] (err3) -| (err_final);
	\end{tikzpicture}
	\caption{Diagrama de Actividades del proceso de emisión y validación del voto.}
	\label{fig:uml_actividades}
	\par\vspace{3pt}{\footnotesize\textit{Fuente: Elaboración propia.}}
\end{figure}

% =============================================================================
